\section{Introduction}\label{sec:introduction}

Let \(A\) be a maximal abelian unital \(*\)-subalgebra, or MASA, of a
complex von Neumann algebra \(M\).
For a finite partition \(\mathcal P=(p_i)_{i=1}^r\) of the unit by
projections in \(A\), write
\[
 C_{\mathcal P}(z)=\sum_{i=1}^r p_i z p_i .
\]
Norm paving asks whether \(C_{\mathcal P}(z)\) can be made close in
operator norm to an element of \(A\). Approximation paving allows a
bounded perturbation of \(z\) within any prescribed strong-operator
neighborhood before choosing the partition.

Kadison and Singer posed the pure-state extension problem for the atomic
diagonal MASA of \(\mathcal B(\ell^2)\) \cite{KS59}. Anderson formulated its
diagonal paving counterpart \cite{Anderson79}, and Marcus, Spielman and
Srivastava resolved it by mixed characteristic polynomials \cite{MSS15}.
For general MASAs, ordinary norm paving of the original operator is too
restrictive: in the separable-predual setting, it characterizes type I
algebras with a normal expectation onto the MASA
\cite[Theorem~3.3]{PV15}. Popa and Vaes introduced strong-operator and
approximation versions of paving and conjectured that every MASA has
both properties \cite[Definitions~2.2--2.3 and Conjecture~2.8(1)]{PV15}.
We resolve the approximation-paving assertion of that conjecture positively.

Popa proved norm paving in ultraproducts of singular finite MASAs,
and hence strong-operator paving in the original finite inclusions
\cite{Popa14}. Popa and Vaes established strong-operator and approximation
paving for type I algebras with separable predual and for Cartan
inclusions in amenable algebras or arising from free ergodic profinite
actions \cite[Sections~3--4]{PV15}. For singular finite MASAs they
obtained quadratic paving bounds \cite[Theorem~5.1]{PV15}, later
sharpened in \cite{PV16}. These restricted cases do not provide the
expectation-free approximation theorem below.

\Needspace{16\baselineskip}
\begin{theorem}\label{thm:main}
There is a universal constant \(C_{\mathrm{ap}}>0\) such that for every
\(0<\varepsilon<1\) there is an integer
\(R(\varepsilon)\leq C_{\mathrm{ap}}\varepsilon^{-6}\)
with the following property.
Let \(M\subseteq\mathcal B(H)\) be a complex von Neumann algebra,
let \(A\subseteq M\) be a MASA, and let \(x=x^*\in M\).
For every finite \(F\subset H\) and every \(\delta>0\), there exist
\(y=y^*\in M\), a partition \(\mathcal P=(p_i)_{i=1}^{R(\varepsilon)}\)
of the unit in \(A\), and \(a=a^*\in A\), such that
\[
 \norm{y}\leq3\norm{x},\qquad
 \norm{(y-x)\xi}<\delta\quad(\xi\in F),
\]
and
\[
 \norm{a}\leq\norm{y},\qquad
 \norm{C_{\mathcal P}(y)-a}\leq\varepsilon\norm{y}.
\]
Zero projections may be included in the partition.
\end{theorem}

Thus the norm-control constant is universal, and the number of
projections is independent of the algebra, its representation, the
operator, and the strong neighborhood. We do not optimize the
sixth-power bound.
The relative error in the theorem is measured against the actual
approximant's norm; a norming test vector and a final clipping argument
will ensure that normalization.

There is also a direct consequence for strong-operator paving when a
normal expectation exists. In its definition, the compression is
estimated on a large projection whose complement lies in a prescribed
strong neighborhood of zero.

\begin{corollary}\label{cor:so-paving}
Suppose that a MASA \(A\subseteq M\) is the range of a normal
conditional expectation. For each \(0<\varepsilon<1\) there is an
integer \(R_s(\varepsilon)\), independent of the inclusion, such that
for every \(x=x^*\in M\) and every strong-operator neighborhood
\(\mathcal V\) of zero, there exist an \(R_s(\varepsilon)\)-partition
\(\mathcal P\) in \(A\), an \(a=a^*\in A\) with
\(\norm{a}\leq\norm{x}\), and a projection \(q\in M\) with
\(\one-q\in\mathcal V\), satisfying
\[
 \norm{q(C_{\mathcal P}(x)-a)q}\leq\varepsilon\norm{x}.
\]
\end{corollary}
The proof is given in Section~\ref{gen:consequences}.

The separate companion \cite[Theorem~1.1]{OAQuadratic26} proves a direct
quadratic strong-operator paving bound for the original operator.
Theorem~\ref{thm:main} instead norm-paves a nearby bounded operator,
with error relative to that approximant's norm. The two main proofs are
independent; the conversion in Corollary~\ref{cor:so-paving} retains its normal-expectation
hypothesis. We claim no quadratic bound for \(R(\varepsilon)\).

\begin{corollary}[Quadratic approximation paving for expected MASAs]
\label{cor:quadratic-expected-approximation}
There is a universal constant \(C>0\) such that for every
\(0<\varepsilon<1\) one can choose an integer
\(N(\varepsilon)\leq C\varepsilon^{-2}\) with the following property.
If the MASA \(A\subseteq M\) is the range of a normal conditional expectation,
then the conclusion of Theorem~\ref{thm:main} holds with
\(N(\varepsilon)\) in place of \(R(\varepsilon)\).
\end{corollary}
The proof is given in Section~\ref{gen:consequences}.

\paragraph{The measured argument.}
The first part of the proof develops a measurable version of
one-sided matrix paving for nonsingular equivalence relations.
The finite paving-polynomial framework of
\cite{RL20,RS21} supplies the starting point. Its stable expected
determinants yield local spectral probabilities and spectral-shift
functions.
A buffered pinning lemma selects a color while increasing the
diagonal penalty by a fixed universal factor and controlling the
increase of positive-root mass.
The initial root estimate from \cite[Lemma~24 and proof of Theorem~25]{RS21}
is derived here from \cite[Theorems~4.1 and~5.1]{MSS15}, using
positive-defect completion and deletion of added Gram columns.

Nonsingularity is handled by passing to the Maharam extension, where
mass transport is invariant, and weighting the matrix by the height
variable. Uniform comparisons between slabs allow the pinning
decisions to descend to a single coloring of the base.
Sparse random updates then control the accumulated cost.
Only local independence is needed; auxiliary graph colorings do not
increase the number of colors in the final paving.
The slab comparison and buffered-pinning estimates isolate the two
uniformity issues in this passage from finite matrices to measurable
operators.

\paragraph{The operator-algebra argument.}
For a MASA in the centralizer of a faithful normal state, we first pave
the algebra generated by its normalizers, using the
Feldman--Moore representation \cite{FM77} and the measured theorem.
The complementary part has diffuse left--right kernels.
We adapt the free-product matrix dilation of
\cite[proof of Theorem~5.1]{PV15}, using random diagonal phases to
obtain asymptotic freeness with respect to the induced state.
The proof controls the moments and their concentration without assuming
traciality, using analytic right multiplication for the modular group.

Finally, a separable reduction preserves the partition bounds.
The corner carrying centralizing normal positive functionals has a
normal expectation; on the complementary corner, almost disjoint
conjugate states produce positive contractions with small diagonal
compressions. A three-corner construction preserves the original
off-diagonal entries until the final partition is chosen.
This is the step that removes the normal-expectation hypothesis.

All new estimates needed for these passages are proved below.
We use established structure theorems for conditional expectations,
standard form, type I algebras and free products, with their hypotheses
specified at the points of application.

Section~\ref{sec:preliminaries} fixes the state-topology tools;
Sections~\ref{poly:section}--\ref{sec:pinning} prove the measured paving
theorem, Section~\ref{sec:expected} handles the state-preserving case,
and Section~\ref{gen:section} removes the remaining hypotheses.
